\documentclass[11pt, a4paper]{article}
\usepackage[a4paper, top=2.5cm, bottom=2.5cm, left=2cm, right=2cm]{geometry}
\usepackage{fontspec}
\usepackage[english, bidi=basic, provide=*]{babel}
\babelprovide[import, onchar=ids fonts]{english}
\babelfont{rm}{TeX Gyre Termes}
\babelfont{sf}{TeX Gyre Heros}
\babelfont{tt}{TeX Gyre Cursor}
\usepackage{enumitem}
\usepackage{amsmath}

\title{Study Notes: The Cell Cycle (Chapter 12)}
\author{Subject: Biology}
\date{}

\begin{document}
\maketitle

\section*{Core Concepts}
\begin{itemize}
    \item \textbf{Purpose of Cell Division:} Reproduction (unicellular), growth, and repair (multicellular).
    \item \textbf{The Goal:} Distribution of identical genetic material (DNA) to two daughter cells.
    \item \textbf{Genome:} A cell's total DNA. Prokaryotes usually have one circular molecule; Eukaryotes have multiple linear molecules.
    \item \textbf{Chromosomes:} Packaged DNA. Somatic cells are diploid ($2n = 46$ in humans); Gametes are haploid ($n = 23$ in humans).
    \item \textbf{Sister Chromatids:} Joined copies of the original chromosome.
\end{itemize}

\section*{Key Definitions}
\begin{itemize}
    \item \textbf{Chromatin:} The complex of DNA and protein that makes up eukaryotic chromosomes.
    \item \textbf{Centromere:} The region where sister chromatids are most closely attached.
    \item \textbf{Mitosis:} Nuclear division resulting in genetically identical nuclei.
    \item \textbf{Cytokinesis:} Division of the cytoplasm.
    \item \textbf{Interphase:} 90\% of the cell cycle; includes G1, S, and G2 phases.
    \item \textbf{Kinetochore:} Protein structure at the centromere for spindle attachment.
\end{itemize}

\section*{Important Formulas / Relations}
\begin{itemize}
    \item \textbf{DNA Length vs. Cell Size:} 2 m of DNA vs. micrometer-scale cell diameter ($\approx 250,000\times$ length).
    \item \textbf{Chromosome Count:} 
    \begin{itemize}
        \item After S phase: Chromosome count remains same, but DNA mass doubles.
        \item After Anaphase start: Chromosome count doubles temporarily as sister chromatids become individual chromosomes.
    \end{itemize}
    \item \textbf{Somatic vs. Gamete:} Somatic = $2n$, Gamete = $n$.
\end{itemize}

\section*{Processes / Workflows}

\subsection*{The Eukaryotic Cell Cycle}
\begin{enumerate}
    \item \textbf{Interphase:}
    \begin{itemize}
        \item \textbf{G1 (First Gap):} Metabolic activity and growth.
        \item \textbf{S (Synthesis):} Metabolic activity, growth, and \textbf{DNA replication}.
        \item \textbf{G2 (Second Gap):} Metabolic activity, growth, and preparation for division.
    \end{itemize}
    \item \textbf{Mitotic (M) Phase:}
    \begin{itemize}
        \item \textbf{Prophase:} Chromosomes condense; spindle starts forming; nucleoli disappear.
        \item \textbf{Prometaphase:} Envelope fragments; spindle attaches to kinetochores.
        \item \textbf{Metaphase:} Chromosomes align at the metaphase plate (center).
        \item \textbf{Anaphase:} \textbf{Separase} cleaves cohesins; chromatids part suddenly; cell elongates.
        \item \textbf{Telophase:} Two daughter nuclei form; nucleoli reappear; chromosomes de-condense.
    \end{itemize}
\end{enumerate}

\subsection*{Cytokinesis: Animal vs. Plant}
\begin{itemize}
    \item \textbf{Animals:} \textbf{Cleavage furrow} forms. A contractile ring of \textbf{actin} and \textbf{myosin} pinches the cell.
    \item \textbf{Plants:} \textbf{Cell plate} forms. Golgi-derived vesicles carry wall materials to the cell center to build a new wall.
\end{itemize}

\subsection*{Prokaryotic Binary Fission}
\begin{itemize}
    \item Chromosome replication starts at the \textbf{origin of replication}.
    \item One origin moves to the opposite pole (actin-like protein involvement).
    \item Cell elongates and membrane pinches inward (tubulin-like protein involvement).
\end{itemize}

\section*{Common Mistakes and Clarifications}
\begin{itemize}
    \item \textbf{Sister Chromatids vs. Chromosomes:} Two sister chromatids make up ONE duplicated chromosome. Once they separate in anaphase, they are called individual chromosomes.
    \item \textbf{Centrosome vs. Centriole:} Centrosome is the region; centrioles are the paired structures inside. Spindles can form without centrioles (plants).
    \item \textbf{Gap Phases:} "Gaps" do not mean the cell is resting; it is metabolically hyperactive.
    \item \textbf{Mitosis vs. Cytokinesis:} Mitosis is nuclear; cytokinesis is cytoplasmic.
\end{itemize}

\section*{Exam-Focused Quick Revision}
\begin{itemize}
    \item \textbf{Checkpoints:}
    \begin{itemize}
        \item \textbf{G1 Checkpoint:} Most critical. "Go-ahead" $\rightarrow$ division; No signal $\rightarrow$ \textbf{G0 phase}.
        \item \textbf{G2 Checkpoint:} Verifies DNA replication is complete.
        \item \textbf{M Checkpoint:} Verifies all kinetochores are attached to spindle fibers.
    \end{itemize}
    \item \textbf{Regulatory Molecules:}
    \begin{itemize}
        \item \textbf{Cyclin:} Concentration fluctuates rhythmically.
        \item \textbf{Cdk (Cyclin-dependent Kinase):} Active only when bound to cyclin.
        \item \textbf{MPF:} Cyclin-Cdk complex that triggers mitosis.
    \end{itemize}
    \item \textbf{Spindle Anatomy:} Centrosomes + Spindle Microtubules + Asters.
\end{itemize}

\end{document}